% RESOURCE-DIRECTSUM REPORT::RESOURCE-GAUGE5::2026-10-18_1520_KST
\documentclass[11pt]{article}
\usepackage{amsmath,amssymb,amsthm,hyperref}
\newtheorem{theorem}{Theorem}
\title{RESOURCE-GAUGE5: A Source-Tailored Locality Escape, Shared-Gauge Direct-Sum Counterexample, and Frozen-Path Selection Hardness}
\date{18 October 2026}
\begin{document}\maketitle
\noindent\textbf{STATUS.} PROVED/EXACT for the stated restricted-model lemmas, pending independent review. Finite experiments are COMPUTATIONAL. RESOURCE-DIRECTSUM1, INDEX-HIDDEN-DIRECTSUM1, CLOCK-INDEPENDENT-GATE1, THRESHOLD-DESCENT1 and P versus NP remain OPEN.

\section{IDEAS PROPOSED BEFORE LITERATURE}
\begin{itemize}
\item[] A. Find a fixed two-probe affine generator admitting an invertible dense output mixer that stays two-probe. This tests the universal-versus-source-tailored quantifier.
\item[] B. Make many individually rank-rich blocks share the same publicly accessible latent vertex state and quantify the failure of additive global rank.
\item[] C. Test pairwise nonlinear postprocessing and an expensive inverse output transform against alternative latent representations.
\item[] D. Determine whether finding an exceptionally small globally favorable cell set is itself difficult when the actual query paths are explicit.
\end{itemize}

\section{LITERATURE CHECK, CORPUS, AND PREVIOUS WORK}
Read project AGENTS.md, CODEX.md, pnp-research governance, root main.tex, DATA-DIRECTSUM1, PCP-GATE1, the recent MCSP cone and clock reports, and locally preserved CODEBASIS-DIRECTSUM2, RESOURCE-ADVICE-LDC3, RESOURCE-READCONGESTION4. The earlier BASIS-COMPLETE1, SELECTOR-ENTROPY1, advice-LDC bounds, and UNIVERSAL-NONLINEAR-LOCALITY-RIGIDITY4 retain their actual quantifiers. The graphic incidence rank formula is standard; see MIT linear-algebra graph-incidence notes. Korten-Pitassi-Impagliazzo, FOCS 2025/ECCC TR25-030, provide a local-PRG/cell-probe connection for a particular hard family, not an arbitrary-generator mixing theorem. Golovnev-Gur-Shinkar, SOSA 2023, give a nonadaptive cell-sampling improvement, not a generic adaptive-program compiler. Clique is NP-complete by Karp (1972). No matching source-tailored dense graph-gauge direct-sum report was found by repository searches for SOURCE-TAILORED-INDEX, READCONGESTION4, or GRAPH-BOUNDARY.

\section{Exact boundary and direct-sum theorem}
Let \(X=(V,E)\) be a loopless undirected graph and \(B\in\mathbb F_2^{E\times V}\) its vertex-edge incidence matrix. For seed \(z\in\mathbb F_2^V\) set \(G_X(z)=Bz\). For an edge subset \(F\), its XOR sum is
\[
  \bigoplus_{e=uv\in F}(z_u\oplus z_v)
  =\bigoplus_{u:\deg_F(u)\equiv1\pmod2}z_u.
\]
Since parity on \(r\) independent bits has exact adaptive deterministic decision-tree query complexity \(r\), a fixed source-tailored linear mixing row is two-seed-probe exactly when its odd-degree boundary has size zero or two, regardless of how many edges the row mixes. Also \(\operatorname{rank}_{\mathbb F_2}(B)=|V|-\operatorname{cc}(X)\), including isolates.

\begin{theorem}[SHARED-GAUGE-DIRECTSUM-NOGO5]
Fix \(v\ge3\) and \(v\) blocks indexed by \(i\in\{0,\ldots,v-1\}\). Block \(i\) contains the \(v-1\) edges \(((i+j)\bmod v,(i+j+1)\bmod v)\), \(0\le j<v-1\), with bit \(h_i(z,j)=z_u\oplus z_w\). Each separate block has rank \(v-1\) and \(2^{v-1}\) possible outputs, but the entire combined \(v(v-1)\)-bit family also has rank \(v-1\), not \(v(v-1)\). There is one uniform public data structure with \(v-1\) variable bits and at most two bit probes per output that recovers all complete blocks.
\end{theorem}
\noindent Proof. Each block is a spanning tree, so its incidence rank is \(v-1\). All blocks operate on the same latent \(v\)-bit seed, with common nullspace the constant vertex assignments, so the joint stacked incidence rank is still \(v-1\). Store \(D_0=0\) publicly and \(D_u=z_u\oplus z_0\) for \(u>0\). Calculate the two edge endpoints from \((i,j)\) in \(O(\log v)\) arithmetic workspace and answer \(D_u\oplus D_w\). The source graph is a uniformly indexed family and requires no per-block endpoint lookup table.

\begin{theorem}[INVERSE-LOCALITY-ILLUSION5]
Let \(m=v-1\), \(b_j=z_0\oplus z_j\), \(c_j=b_j\oplus b_{j+1}\) for \(1\le j<m\), and \(c_m=b_m\). The forward transformation \(b\mapsto c\) is invertible and has at most two input-symbol probes per output. The inverse is \(b_j=\bigoplus_{\ell=j}^m c_\ell\), requiring exactly \(m-j+1\) probes if restricted to raw \(c\)-symbol cells. Nevertheless the target \(c\) admits an alternative \(m\)-bit gauge representation answering every target query in at most two probes.
\end{theorem}
\noindent Proof. The forward matrix is upper bidiagonal with diagonal one, and its inverse is an all-ones suffix matrix. Parity is sensitive to every suffix bit. But retaining the \(b\)-gauge table directly answers \(c_j\) by one or two bit reads. Thus inverse difficulty under a forced representation is not a lower bound over all representations.

\begin{theorem}[PAIRWISE-FACTOR-AND-RANGE-NOGO5]
For arbitrary public uniformly computable functions \(\varphi_{i,j}:\{0,1\}^2\to\{0,1\}\) on these edge endpoints, retaining the \(v\) raw vertex bits yields at most two probes per target query, even if the output function is nonlinear. The smaller \(v-1\) gauge need not suffice: all-zero and all-one vertex seeds have the same gauge but different edge-AND outputs. Furthermore membership in the image of any explicit XOR-labeled graph-edge generator, and a compatible seed if one exists, are decidable in \(O(|V|+|E|)\) graph operations.
\end{theorem}
\noindent Proof. The general nonlinear table reads its two named raw seed bits and evaluates the public gate. For affine XOR outputs, choose a root potential in each connected component, propagate edge parities through a spanning forest, and check all remaining edges. An inconsistent nonforest edge closes a contradictory cycle; otherwise propagated vertex potentials realize the labels. Hence a simple graphic XOR generator is not a hard range-avoidance seed by itself.

\section{A charged-discovery barrier}
\begin{theorem}[FROZEN-TWO-PROBE-SELECT-NPC5]
Given an explicit family of two-cell query paths \(P_e\subseteq[S]\), a cell budget \(\Delta\) and target \(r\), deciding whether some \(C\subseteq[S]\), \(|C|\le\Delta\), contains at least \(r\) paths \(P_e\subseteq C\), is NP-complete, already for fixed nonadaptive paths.
\end{theorem}
\noindent Proof. Membership in NP is immediate. Reduce a graph's CLIQUE instance \((G,k)\) by creating \(P_e=\{u,v\}\) for every edge \(e=uv\), taking \(\Delta=k\) and \(r={k\choose2}\). A set of at most \(k\) vertices induces at most \({k\choose2}\) edges, with equality iff it is a \(k\)-clique. This is the usual densest-\(k\)-subgraph endpoint. The theorem concerns optimization of frozen replay-path certificates, NOT an optimal semantic data structure, nor unconditional computational lower bounds for KPI strings. A program may use other query paths or precomputed aggregates.

\section{Reproducible finite audit}
Standalone Python source: resource_directsum5_test.py (SHA-256 52d44b7c21bd8f01d42a69916629854301f692e9b56158f86bc69ab9accaa055), raw results: RESOURCE_GAUGE5_20261018_raw.json. One deterministic, standard-library batch took about 2.2 s; no random seeds or asymptotic extrapolation. At \(v=6\): six rank-five blocks, sum rank 30, joint rank five, 64 seeds, 1,920 XOR checks, 7,680 nonlinear pairwise checks, 4,096 boundary identities, and 1,920 one-bit off-image mutations passed. At \(v=8\): eight rank-seven blocks, sum rank 56, joint rank seven, 256 seeds, 14,336 XOR checks, 57,344 nonlinear controls, 65,536 boundary identities, and 14,336 single-output-bit off-image mutations passed. All 1,024 simple graphs on five vertices and 32,768 on six vertices passed 167,936 exact CLIQUE-to-selector equivalence checks.

\section{Status, duplication guard, next target}
PROVED/EXACT (restricted): graphic parity boundary criterion, uniform shared-gauge complete-block counterexample, forced-inverse versus alternate-representation separation, pairwise public nonlinear factorization, affine range membership, and frozen replay-path selection NP-completeness. COMPUTATIONAL: finite test counts. Previously established results are retained: BASIS-COMPLETE1, SELECTOR-ENTROPY1, READ-CONGESTION-FULLBLOCK4, UNIVERSAL-NONLINEAR-LOCALITY-RIGIDITY4, resource-aware LDC bounds. No theorem here applies to a KPI remote point outside the easy graph-code image. No automatic cell-probe-to-RAM transfer or MCSP slope \(\rho>\lambda\) is claimed.

Highest-value next lemma: SOURCE-TAILORED-NONGAUGE6. Start with the precise published KPI source, prove it lacks a cheap globally shared latent factorization while its off-range hard-side construction remains semantically gated, and test adaptive index/address selection against arbitrary alternative memory encodings. Any claimed contradiction must remain in the same exact probe, encoding, and succinctness model. P versus NP remains OPEN.

\begin{thebibliography}{9}
\bibitem{kpi} O. Korten, T. Pitassi, R. Impagliazzo, \emph{Stronger Cell Probe Lower Bounds via Local PRGs}, FOCS 2025, ECCC TR25-030. \url{https://eccc.weizmann.ac.il/report/2025/030/}.
\bibitem{ggs} A. Golovnev, T. Gur, I. Shinkar, \emph{Derandomization of Cell Sampling}, SOSA 2023. \url{https://doi.org/10.1137/1.9781611977585.ch26}.
\bibitem{karp} R. Karp, \emph{Reducibility Among Combinatorial Problems}, 1972. \url{https://doi.org/10.1007/978-1-4684-2001-2_9}.
\end{thebibliography}
\end{document}
% END_OF_VERIFIED_REPORT::RESOURCE-DIRECTSUM::RESOURCE-GAUGE5::P_VS_NP_OPEN
